\chapter{Contribuțiile proprii}

Capitolul inițial subliniază contribuțiile proprii la disertație, concentrându-se pe elementele originale de proiectare, integrare, adaptare și validare experimentală. Nu este scopul lucrării să dezvolte o nouă teorie a rețelelor recurente; mai degrabă, este de a crea un cadru aplicativ și metodologic coerent pentru a examina în mod sistematic impactul valorilor lipsă asupra predicției seriilor de timp. Din această perspectivă, arhitectura experimentală dezvoltată și procesele de preprocesare, imputare, antrenare, evaluare și interpretare a rezultatelor sunt dovezi ale contribuției personale.

Prin utilizarea acestei metode, lucrarea depășește o simplă prezentare teoretică a unor metode cunoscute și propune un cadru coerent de analiză comparativă care poate fi utilizat pentru a examina relația dintre calitatea datelor de intrare și performanța modelelor predictive. În continuare, sunt prezentate contribuțiile proprii prin principalele direcții de dezvoltare și adaptare realizate în cercetare.

\section{Proiectarea unui cadru experimental unitar}

Proiectarea unui cadru experimental unitar pentru a examina efectul valorilor lipsă asupra seriilor de timp este o primă contribuție proprie. Acest cadru combină etapele de încărcare a datelor, injectarea controlată a valorilor lipsă, implementarea metodelor de imputare, configurarea modelelor recurente, pregătirea, generarea de predicții și evaluarea prin metrici standardizate într-un singur flux de lucru. Acest aport este unic pentru că toate aceste etape au fost integrate într-un sistem coerent, reproductibil și ușor de extins pentru experimente suplimentare.

Soluția sugerată permite analiza comparativă a întregului lanț experimental în condiții controlate. Acest lucru este diferit de abordările fragmentate, în care fiecare etapă este tratată separat. Prin urmare, impactul fiecărei decizii metodologice poate fi văzut mai clar, iar rezultatele pot fi interpretate mai atent și mai relevante pentru scopul general al lucrării.

\section{Adaptarea mecanismelor de tratare a valorilor lipsă}

Abordarea și adaptarea mai multor abordări pentru tratarea valorilor lipsă într-un mediu experimental comun este o contribuție semnificativă a lucrării. Aplicația creată a folosit atât metode simple, cum ar fi propagarea valorii ultime observate, completarea cu zero sau cu media coloanei, cât și metode mai avansate, cum ar fi imputarea predictivă bazată pe \textit{Random Forest}. În ciuda faptului că aceste metode sunt prezente în literatura și practica analizei datelor, avantajul lor specific constă în adaptarea lor la același mediu experimental și în evaluarea lor sistematică asupra acelorași serii de timp și arhitecturi de predicție.

În urma acestei integrări, a fost posibil să se efectueze o analiză detaliată a relației dintre modelul de forecasting utilizat și metoda de imputare. Ca urmare, nu numai că se utilizează tehnici cunoscute, ci se compară și organizează acestea într-un cadru care subliniază impactul lor real asupra performanței predictive.

\section{Implementarea și compararea arhitecturilor recurente}

O altă contribuție proprie este implementarea și compararea a trei arhitecturi recurente (RNN, LSTM și GRU) care sunt relevante pentru predicția seriilor de timp în aceleași aplicații. Aceste modele pot fi evaluate în condiții similare folosind seturi de date identice, strategii de împărțire a seriilor identice și un cadru general de evaluare identic. Această metodă a permis cercetarea diferențelor de comportament dintre arhitecturi în ceea ce privește datele incomplete și tehnicile de completare utilizate.

Calitatea unică a acestei contribuții provine din faptul că analiza nu se limitează la performanța unui singur model; mai degrabă, se concentrează pe modul în care arhitectura, calitatea datelor și strategia de imputare se influențează reciproc. Lucrarea oferă, prin urmare, o perspectivă comparativă cu valoare practică, care este utilă în procesul de a alege cea mai bună soluție în funcție de contextul datelor care sunt disponibile.

\section{Extinderea analizei prin studii de configurare}

În afară de comparația de bază dintre modele și metode de imputare, lucrarea aduce o contribuție suplimentară prin extinderea evaluării experimentale. Lungimea secvenței, dimensiunea stratului ascuns, numărul de straturi, rata de învățare, scăderea, optimizatorul și funcția de pierdere sunt hiperparametri importanți pentru care au fost integrate mecanisme în aplicația dezvoltată pentru a analiza efectele normalizării datelor. Această extindere a metodologiei permite observarea sensibilității modelelor la diferite configurații de antrenare. De asemenea, ajută la susținerea interpretării mai detaliate a rezultatelor.

Prin includerea acestor studii, lucrarea depășește nivelul implementării demonstrative și devine o platformă de analiză experimentală. În acest sens, contribuția proprie constă în transformarea aplicației de predicție într-un instrument capabil să analizeze impactul valorilor inexistente în raport cu mai multe elemente metodologice.

\section{Integrarea învățării federate în aplicația experimentală}

O contribuție proprie suplimentară este integrarea unui flux de învățare federată în aceeași aplicație experimentală. În locul antrenării unui singur model pe toate datele, seria este împărțită pe utilizatori, iar modelele sunt antrenate local înainte de agregarea rezultatelor într-un model global. Această structură permite evaluarea comparativă a performanței pentru fiecare utilizator și a efectului agregării la nivel central.

Această integrare este importantă deoarece leagă direct teoria învățării federate de o implementare practică folosită pentru seriile de timp. În loc să trateze federated learning ca pe o idee separată, lucrarea îl folosește ca instrument pentru a observa cum se modifică predicția atunci când datele sunt distribuite și când actualizarea globală depinde de contribuția fiecărui participant~\cite{mcmahan2017fedavg, kairouz2021fl}.

\section{Integrarea experimentelor de ansamblu}

De asemenea, dezvoltarea și compararea experimentelor de ansamblu reprezintă o contribuție proprie distinctă. Aplicația include atât o topologie independentă, în care mai multe modele contribuie separat la evaluare, cât și o topologie secvențială bazată pe reziduuri. Această dublă abordare permite urmărirea modului în care diferite scheme de combinare influențează metricile finale și stabilitatea predicției.

Contribuția nu constă doar în rularea unor metode cunoscute, ci în integrarea lor în același cadru comparativ în care sunt deja studiate valorile lipsă, imputarea și modelele recurente. Prin această extindere, lucrarea arată cum poate fi folosit ansamblul nu doar ca o tehnică izolată, ci ca parte dintr-o metodologie mai amplă de evaluare a performanței predictive~\cite{dietterich2000ensemble, zhou2012ensemble}.

\section{Dezvoltarea unei interfețe interactive pentru experimentare}

Dezvoltarea unei interfețe interactive care permite configurarea dinamică a experimentelor este o componentă originală semnificativă a lucrării. Utilizatorul poate alege setul de date, procentul de valori lipsă, metoda de imputare, arhitectura modelului, modul de normalizare și hiperparametrii de antrenare prin intermediul aplicației realizate. Apoi poate vizualiza rezultatele într-o formă vizuală și comparativă. Cadrul teoretic și metodologic este transformat într-un instrument util pentru testare și explorare prin intermediul acestei interfețe.

În acest caz, contribuția proprie implică crearea unei soluții care permite reproducerea experimentelor și efectuarea unei analize interactive a rezultatelor. Acest aspect este util atât pentru validarea rezultatelor lucrării, cât și pentru utilizarea ulterioară în scop educațional sau aplicativ. Din această perspectivă, formularea conceptuală a problemei nu este singura care oferă valoare lucrării; mai degrabă, există o implementare funcțională.

\section{Validarea pe seturi de date distincte}

Alegerea a două seturi de date cu proprietăți distincte, DSMTS și \textit{Bitcoin Historical Data}, susțin caracterul aplicativ al contribuțiilor proprii. Cadrul sugerat a fost testat atât într-un mediu mai controlat, cât și într-unul cu volatilitate ridicată și dinamică temporală complexă cu ajutorul acestor seturi. Deci, concluziile nu se bazează pe un caz specific; mai degrabă, ele se bazează pe rezultatele diferitelor condiții experimentale.

Această decizie metodologică este utilă, deoarece oferă o bază mai solidă pentru interpretarea performanțelor observate și pentru evaluarea credibilității soluției propuse. În plus, utilizarea diferitelor seturi de date extinde relevanța practică a lucrării și extinde domeniul de aplicare a rezultatelor sale.

\section{Sinteza contribuțiilor originale}

În concluzie, există o serie de linii complementare care pot fi utilizate pentru a formula contribuțiile individuale ale disertației: dezvoltarea unui cadru experimental unitar pentru analiza valorilor lipsă în seriile de timp; adaptarea și compararea sistematică a diferitelor strategii de imputare; implementarea arhitecturilor RNN, LSTM și GRU într-un mod comparativ; extinderea studiului prin analize de hiperparametri și normalizare; integrarea învățării federate și a experimentelor de ansamblu în aceeași aplicație; crearea unei interfețe interactive pentru configurarea experimentelor și validarea metodologiei pe diferite seturi de date. Aceste componente stabilesc componenta originală a lucrării și îi dau consistență din punct de vedere metodologic și aplicativ.

Scopul disertației este de a investiga modul în care metode consacrate au fost combinate, adaptate și utilizate comparativ într-un context unitar. Acest lucru se datorează faptului că disertația se concentrează în mod specific pe modul în care impactul valorilor absente asupra performanței modelelor de forecasting. Principalul aport personal al lucrării este această perspectivă aplicativă și sistematică.
